% Standalone problem document, generated from the adjacent README.md.
% Compile with XeLaTeX. Mathematical content is maintained in Markdown.
\documentclass[11pt,a4paper]{article}
\usepackage[fontsize=10.5pt]{scrextend}
\usepackage[margin=22mm,headheight=15pt,headsep=9mm,footskip=11mm]{geometry}
\usepackage{fontspec}
\setmainfont{texgyrepagella}[Extension=.otf,UprightFont=*-regular,BoldFont=*-bold,ItalicFont=*-italic,BoldItalicFont=*-bolditalic]
\setsansfont{texgyreheros}[Extension=.otf,UprightFont=*-regular,BoldFont=*-bold,ItalicFont=*-italic,BoldItalicFont=*-bolditalic]
\setmonofont{texgyrecursor}[Extension=.otf,UprightFont=*-regular,BoldFont=*-bold,ItalicFont=*-italic,BoldItalicFont=*-bolditalic,Scale=MatchLowercase]
\usepackage{amsmath,amssymb}
\usepackage{unicode-math}
\setmathfont{texgyrepagella-math.otf}
\usepackage{xcolor,microtype,enumitem,needspace,fancyhdr}
\usepackage{bookmark}
\definecolor{ink}{HTML}{17344B}
\definecolor{accent}{HTML}{197D80}
\definecolor{muted}{HTML}{596773}
\hypersetup{colorlinks=true,linkcolor=accent,urlcolor=accent,pdftitle={IE-17 - Monotonic
optimal backward error along
LSMR},pdfauthor={Open Problems in Numerical Linear Algebra}}
\urlstyle{same}
\setlength{\parindent}{0pt}
\setlength{\parskip}{5pt plus 1pt minus 1pt}
\linespread{1.04}
\setlength{\emergencystretch}{3em}
\widowpenalty=10000
\clubpenalty=10000
\displaywidowpenalty=10000
\allowdisplaybreaks[1]
\setlist{nosep,leftmargin=1.5em}
\providecommand{\tightlist}{\setlength{\itemsep}{0pt}\setlength{\parskip}{0pt}}
\providecommand{\pandocbounded}[1]{#1}
\pagestyle{fancy}
\fancyhf{}
\fancyhead[L]{\sffamily\footnotesize\color{muted}OPEN PROBLEMS IN NUMERICAL LINEAR ALGEBRA}
\fancyhead[R]{\sffamily\bfseries\footnotesize\color{accent}IE-17}
\fancyfoot[L]{\parbox[b]{0.9\textwidth}{\sffamily\fontsize{8}{10}\selectfont\color{muted}Editorial ratings; a bounded search cannot certify that a problem remains unsolved.\\Literature check: 2026-09-14}}
\fancyfoot[R]{\sffamily\footnotesize\color{muted}\thepage}
\renewcommand{\headrulewidth}{0.3pt}
\renewcommand{\headrule}{\hbox to\headwidth{\color{accent}\leaders\hrule height \headrulewidth\hfill}}
\makeatletter
\renewcommand{\section}{\Needspace{5\baselineskip}\@startsection{section}{1}{0pt}{12pt plus 2pt minus 2pt}{4pt}{\sffamily\bfseries\large\color{ink}}}
\renewcommand{\subsection}{\Needspace{4\baselineskip}\@startsection{subsection}{2}{0pt}{9pt plus 1pt minus 1pt}{3pt}{\sffamily\bfseries\normalsize\color{ink}}}
\makeatother
\setcounter{secnumdepth}{0}
\begin{document}
{\sffamily\small\color{muted}Linear systems and elimination\par}
\vspace{3pt}
{\raggedright\hyphenpenalty=10000\exhyphenpenalty=10000\sffamily\bfseries\fontsize{21}{24}\selectfont\color{ink}Monotonic
optimal backward error along LSMR\par}
\vspace{7pt}
{\sffamily\small\textbf{Difficulty:} challenging\quad\textbf{Importance:} interesting
to the community\par}
{\sffamily\small\bfseries\color{accent}Status: Lean verified\par}
\vspace{4pt}
{\color{accent}\hrule height 0.8pt}
\vspace{6pt}
\subsection{Independently reviewed resolution -
2026-09-11}\label{independently-reviewed-resolution---2026-09-11}

\textbf{Negative resolution.} Sections 1-4 give one exact
full-column-rank \(4\times3\) LSMR example for which both displayed
errors increase from the first to the second nonzero iterate. The
matrix-only spectral backward error satisfies
\(\mu(x_1)^2\le1979/2000<99/100<\mu(x_2)^2\), and the specified
approximation also strictly increases. The right-hand side stays fixed.
This settles the canonical spectral-norm formulation; the cited SISC
paper uses a different default norm convention, so no Frobenius-error
conclusion is inferred.

\textbf{Author:} Matthew J. Colbrook, Department of Applied Mathematics
and Theoretical Physics, University of Cambridge.
\href{https://github.com/ajt60gaibb/OpenProblemsInNLA/blob/main/references/colbrook-recovered-2026-09-11/manuscripts/IE-17.pdf}{Complete
manuscript},
\href{https://github.com/ajt60gaibb/OpenProblemsInNLA/blob/main/references/colbrook-recovered-2026-09-11/verification/reviews/IE-17-review.md}{independent
proof review}, and
\href{https://github.com/ajt60gaibb/OpenProblemsInNLA/blob/main/references/colbrook-recovered-2026-09-11/README.md}{submission
record}. The supplied notes were reconstructed with substantial AI
assistance. This is independent agent verification, not external human
peer review or formal proof-assistant certification; no novelty or
priority claim is made.

The difficulty, importance and rating rationale below are historical
assessments of the original open target. Original statements, references
and dated audits are preserved.

\textbf{Rating rationale:} Challenging reflects monotonicity of an
optimization-defined error along coupled Krylov iterates; community
impact is a stopping and reliability guarantee for a widely used
least-squares method.

\subsection{Lean proof and verification evidence ---
2026-09-14}\label{lean-proof-and-verification-evidence-2026-09-14}

\textbf{Both original monotonicity claims have Lean-verified negative
answers.} The complete pair of monotonicity claims is disproved by the
exact first two nonzero iterates of undamped, zero-initial-guess LSMR on
a real 4-by-3 system. The formalization proves the genuine
minimum-length normal-residual minimizers over the Krylov spaces, the
matrix-only spectral backward-error infimum and attainment, and the
prescribed Moore--Penrose projector approximation. Both errors strictly
increase; the right-hand side remains fixed.

\href{https://github.com/sidneyholden1/OpenProblemsInNLA/tree/d45afc8a197ffeeff94abfca59ef93acda22efa5/linear-systems-and-elimination/IE-17/lean}{Immutable
formal proof}. The checked exports are \textit{NLA.\allowbreak{}IE17.\allowbreak{}iterates},
\textit{NLA.\allowbreak{}IE17.\allowbreak{}backward\_\allowbreak{}increase},
\textit{NLA.\allowbreak{}IE17.\allowbreak{}approximation\_\allowbreak{}increase},
\textit{NLA.\allowbreak{}IE17.\allowbreak{}counterexample}. Lean 4.33.1 uses pinned Mathlib
\textit{0df444a3} and LeanCert \textit{621a43d7};
\href{https://github.com/sidneyholden1/OpenProblemsInNLA/tree/codex/lean-ie17/linear-systems-and-elimination/IE-17/lean/README.md}{complete
pins and reproduction instructions} are retained. All four exported
transitive axiom closures contain only \textit{propext},
\textit{Classical.choice} and \textit{Quot.sound}.

\textbf{Formalization:} Sidney Holden, Center for Computational Biology,
Flatiron Institute, Simons Foundation, with OpenAI Codex assistance.
Matthew J. Colbrook retains mathematical authorship. Two independent
statement reviews preceded implementation, and two independent final
proof reviews passed. Kernel-only LeanCert and
\href{https://github.com/sidneyholden1/OpenProblemsInNLA/actions/runs/34862174739}{fresh
isolated Linux Comparator and kernel verification} passed, including
rejection controls.
\href{https://github.com/sidneyholden1/OpenProblemsInNLA/tree/codex/lean-ie17/linear-systems-and-elimination/IE-17/lean/verification/linux-2026-09-14/OPERATIONAL-REVIEW.md}{Original
artifact, exact checked inputs and independent operational audit} are
retained. This is AI-assisted formal verification, not external human
peer review or source-author endorsement.

\newpage

Let \(A\in\mathbb R^{m\times n}\) and \(b\in\mathbb R^m\). In exact
arithmetic, start LSMR at \(x_0=0\): equivalently, \(x_k\) minimizes
\(\|A^T(b-Ax)\|_2\) over \(\mathcal K_k(A^TA,A^Tb)\). Work up to its
exact termination and use its minimum-length iterate if necessary. Let
\(r=b-Ax\) and define the matrix-only normwise backward error

\[\mu(x)=\min\{\|E\|_2:(A+E)^T((A+E)x-b)=0\}.\]

For \(x\ne0\), put \(K_x=[A^T,\,(\|r\|_2/\|x\|_2)I]^T\),
\(v_x=[r^T,0^T]^T\), and

\[\widetilde\mu(x)=\frac{\|K_xK_x^{\dagger}v_x\|_2}{\|x\|_2}.\]

Here \(\dagger\) is the Moore--Penrose inverse; at an exact
least-squares solution set both errors to zero.

Are both sequences \(\mu(x_k)\) and \(\widetilde\mu(x_k)\)
nonincreasing, for successive nonzero LSMR iterates? These two closely
related claims are counted together, as in the source. The right-hand
side \(b\) is kept fixed in the backward-error model.

A positive answer would justify backward-error stopping decisions
without a later iteration making the current iterate less backward
accurate. The known monotonicity of \(\|r_k\|_2\) and \(\|A^Tr_k\|_2\)
does not establish this statement.

\subsection{References}\label{references}

D. C.-L. Fong, \emph{Minimum-Residual Methods for Sparse Least-Squares
Using Golub--Kahan Bidiagonalization}, Stanford dissertation 2011,
§7.2.1 and definitions (4.1),(4.5), printed pp.~56--57,118
(\href{https://web.stanford.edu/group/SOL/dissertations/david-fong-thesis-online.pdf}{primary
PDF}). Fong and M. Saunders, \emph{LSMR: An iterative algorithm for
sparse least-squares problems}, SISC 33(2011),2950--2971, §§6.1--6.2,
Fig7.5
(\href{https://web.stanford.edu/group/SOL/software/lsmr/LSMR-SISC-2011.pdf}{author
PDF}). E. Hallman and M. Gu, \emph{LSMB: Minimizing the backward error
for least-squares problems}, SIMAX 39(2018),1295--1317
(\href{https://erhallma.math.ncsu.edu/papers/hallman2018lsmb.pdf}{author
PDF}).

\subsection{Status check --- 2026-09-08}\label{status-check-2026-09-08}

Exact-title and ``LSMR optimal backward error monotonicity
conjecture/proof/counterexample'' searches found no resolution. The
LSMB2018 discussion of nonmonotonicity concerns LSQR, and its proposed
method differs from LSMR. Hallman's May 2026 arXiv:2605.09211 develops
backward-error formulas and estimates, without resolving this
iterate-monotonicity conjecture. No recent explicit reaffirmation was
located.

\subsection{Audit update --- 2026-09-10}\label{audit-update-2026-09-10}

Rechecked Fong's
\href{https://web.stanford.edu/group/SOL/dissertations/david-fong-thesis-online.pdf}{dissertation},
§7.2.1, printed p.~118, which separately conjectures monotonicity of
both displayed errors. \href{https://arxiv.org/abs/2605.09211}{Hallman
(2026)} studies a variational equation and lower bound for the backward
error; it does not assert this iterate monotonicity. Searches for LSMR
optimal-backward-error monotonicity located no proof or counterexample.
\end{document}
